% =========================================================================
\section{Estimación del error al truncar el desarrollo}\label{sec:error}
% =========================================================================
La pregunta obligada ahora es: ¿cuántos términos del desarrollo debemos considerar? Todo depende del error con el que queramos determinar el
potencial. Pero conocer el error exige conocer el potencial exacto, que es precisamente lo que no sabemos calcular y la razón por la que hacemos el
desarrollo. Empezamos, por eso, con una situación en la que sí lo conocemos, y después veremos que, bajo una hipótesis muy natural, el error puede
acotarse sin conocer el potencial exacto.

\begin{definicion}[Error relativo de truncamiento]\label{def:error}
Sea $\phi^{(N)}=\sum_{\ell=0}^{N}\phimul{\ell}$ el desarrollo multipolar truncado en el orden $N$. Su \emph{error relativo} en el punto $\vc{r}$ es
\begin{equation}\label{eq:error}
  e_{N}(\vc{r})=\left\lvert\frac{\phi_{\mathrm{exacto}}(\vc{r})-\phi^{(N)}(\vc{r})}{\phi_{\mathrm{exacto}}(\vc{r})}\right\rvert\times100\ (\%).
\end{equation}
\end{definicion}

% -------------------------------------------------------------------------
\subsection{Un caso con solución exacta: la carga puntual}\label{sec:error-carga}
% -------------------------------------------------------------------------
\begin{ejemplo}[Carga puntual en el origen]\label{ej:carga-origen}
Si la carga $q$ está en el origen, $\phi_{\mathrm{exacto}}=q/(4\pi\eo r)$ coincide con el término monopolar: todos los demás momentos son nulos ($\vp=q\,\vc{0}=\vc{0}$,
$\Qjk=q\,(0-0)=0$, \dots) y, tomando un solo término, el error es del $0\%$.
\end{ejemplo}

\begin{ejemplo}[Carga puntual fuera del origen]\label{ej:carga-fuera}
Sea ahora la carga $q$ en $\rp=a\,\uz$ (\cref{fig:error-carga}). El potencial exacto es
\[
  \phi_{\mathrm{exacto}}(\vc{r})=\frac{1}{4\pi\eo}\,\frac{q}{R}=\frac{q}{4\pi\eo\sqrt{r^{2}+a^{2}-2ra\cos\theta}},
\]
y, por el \cref{teo:1R-legendre} con $r'=a$ y $\theta'=\theta$, su desarrollo multipolar es $\phi=\frac{q}{4\pi\eo}\sum_{\ell}\frac{a^{\ell}}{r^{\ell+1}}\Pl{\ell}(\cos\theta)$. Los
dos primeros términos son
\[
  \phimul{0}=\frac{q}{4\pi\eo r},\qquad \phimul{1}=\frac{\vp\cdot\ur}{4\pi\eo r^{2}}=\frac{q\,\rp\cdot\ur}{4\pi\eo r^{2}}=\frac{q\,a\cos\theta}{4\pi\eo r^{2}},
  \qquad \phi^{(1)}=\frac{q}{4\pi\eo}\Bigl(\frac1r+\frac{a\cos\theta}{r^{2}}\Bigr).
\]
La \cref{tab:error-carga} recoge el error cometido, con $q=1\,\mathrm{nC}$ y $a=1\,\mathrm{m}$, para $\theta=\pi/2$ y $\theta=0$.
\end{ejemplo}

\begin{table}[ht]
\centering
\caption{Potencial exacto y aproximado (términos monopolar y dipolar) de una carga $q=1\,\mathrm{nC}$ situada en $a\,\uz$ con $a=1\,\mathrm{m}$, y error relativo
$e_{1}$ (\cref{def:error}). Se ha usado $1/(4\pi\eo)=8{,}98755\cdot10^{9}\,\mathrm{N\,m^{2}/C^{2}}$.}
\label{tab:error-carga}
\renewcommand{\arraystretch}{1.25}
\small\setlength{\tabcolsep}{3.5pt}
\begin{tabular}{@{}c@{\quad}ccc@{\quad}ccc@{}}
  \toprule
  & \multicolumn{3}{c}{$\theta=\pi/2$} & \multicolumn{3}{c}{$\theta=0$}\\
  \cmidrule(r{0.8em}){2-4}\cmidrule{5-7}
  $r$ & $\phi_{\mathrm{exacto}}$ (V) & $\phi^{(1)}$ (V) & $e_{1}$ (\%) & $\phi_{\mathrm{exacto}}$ (V) & $\phi^{(1)}$ (V) & $e_{1}$ (\%)\\
  \midrule
  $10\,a$   & $0{,}894295$  & $0{,}898755$  & $0{,}499$              & $0{,}998617$  & $0{,}988631$  & $1{,}00$\\
  $100\,a$  & $0{,}0898710$ & $0{,}0898755$ & $5{,}00\cdot10^{-3}$   & $0{,}0907834$ & $0{,}0907743$ & $1{,}00\cdot10^{-2}$\\
  $1000\,a$ & $8{,}987547\cdot10^{-3}$ & $8{,}987552\cdot10^{-3}$ & $5{,}00\cdot10^{-5}$ & $8{,}996548\cdot10^{-3}$ & $8{,}996539\cdot10^{-3}$ & $1{,}00\cdot10^{-4}$\\
  \bottomrule
\end{tabular}
\end{table}

% =========================================================================
% Figura: error relativo e_1 (términos monopolar + dipolar) del potencial de
% una carga q en a u_z, en función de r/a, para theta = 0, pi/4, pi/2 (escala
% logarítmica). Datos exactos: figuras/datos/error_carga.dat.
% =========================================================================
\begin{figure}[ht]
\centering
\begin{minipage}[c]{0.3\textwidth}\centering
\begin{tikzpicture}[x=0.88cm, y=0.88cm, >=Stealth]
  \draw[eje] (0,-0.3) -- (0,2.6) node[above, etiq] {$z$};
  \draw[eje] (-0.3,0) -- (2.9,0) node[right, etiq] {$y$};
  \draw[eje] (0,0) -- (-0.8,-0.7) node[below left=-2pt, etiq] {$x$};
  \coordinate (O) at (0,0);
  \coordinate (Q) at (0,1.0);
  \coordinate (P) at (2.5,1.7);
  \draw[vect, rojoteo] (O) -- (Q) node[pos=0.5, left=1pt, etiq, text=rojoteo] {$\rp=a\,\uz$};
  \draw[vect] (O) -- (P) node[pos=0.55, below right=-2pt, etiq] {$\vc{r}$};
  \draw[->, line width=0.6pt, grispizarra, densely dashed] (Q) -- (P) node[pos=0.5, above=1pt, etiq] {$\vc{R}$};
  \pic[draw=magacero, line width=0.5pt, angle radius=0.55cm, "$\theta$", angle eccentricity=1.45, font=\small] {angle=P--O--Q};
  \node[cargapos] at (Q) {$+$};
  \fill[magprof] (P) circle (1.5pt) node[right=2pt, etiq] {$P$};
\end{tikzpicture}
\end{minipage}\hfill
\begin{minipage}[c]{0.68\textwidth}\centering
\begin{tikzpicture}
\begin{loglogaxis}[marcoMMA, width=8.2cm, height=5.6cm, xmin=1.5, xmax=1000, ymin=1e-5, ymax=60,
                   xlabel={$r/a$}, ylabel={$e_{1}$ (\%)}, every axis plot/.style={line width=1.1pt},
                   log basis y=10, ytick={1e-4,1e-2,1,10}, xtick={2,10,100,1000},
                   xticklabels={$2$,$10$,$100$,$1000$},
                   legend pos=north east, legend cell align=left]
  \addplot[mmaverde!85!black] table[x=x, y=e0]  {figuras/datos/error_carga.dat};
  \addplot[mmanaranja]        table[x=x, y=e45] {figuras/datos/error_carga.dat};
  \addplot[mmaazul]           table[x=x, y=e90] {figuras/datos/error_carga.dat};
  \legend{$\theta=0$, $\theta=\pi/4$, $\theta=\pi/2$}
\end{loglogaxis}
\end{tikzpicture}
\end{minipage}
\caption{Error relativo $e_{1}$ al aproximar el potencial de una carga puntual situada en $a\,\uz$ por sus términos monopolar y dipolar (\cref{ej:carga-fuera}). En escala
doble logarítmica las curvas son rectas de pendiente $-2$: el error decrece como $(a/r)^{2}$, con un coeficiente $\abs{\Pl{2}(\cos\theta)}$ que vale $1$ sobre el eje, $1/4$
para $\theta=\pi/4$ y $1/2$ en el plano ecuatorial. La curva $\theta=0$ coincide con la cota $100\,(a/r)^{2}$ del primer término despreciado, y la caída brusca de la curva $\theta=\pi/4$ cerca de $r=2a$ corresponde a un cambio de signo del error. Los valores para $r=10a,100a,1000a$ son los de la
\cref{tab:error-carga}.}
\label{fig:error-carga}
\end{figure}


\begin{observacion}[Lectura de los resultados]\label{obs:error-carga}
Al multiplicar la distancia por diez, el error se divide por cien: decrece como $(a/r)^{2}$, que es el peso relativo del primer término despreciado, el
cuadrupolar ($\phimul{2}/\phimul{0}=(a/r)^{2}\Pl{2}(\cos\theta)$). El coeficiente es $\abs{\Pl{2}(\cos\theta)}$: vale $1$ sobre el eje y $\frac12$ en el plano ecuatorial, de
ahí que el error sobre el eje sea el doble (\cref{fig:error-carga}). En la dirección del cono nodal, $\cos^{2}\theta=1/3$, el término cuadrupolar se anula y el error cae
un orden más, al $(a/r)^{3}$.
\end{observacion}

% -------------------------------------------------------------------------
\subsection{Cota del error para distribuciones de un solo signo}\label{sec:cota}
% -------------------------------------------------------------------------
Existe una situación en la que, aunque no conozcamos el potencial exacto, podemos estimar el error cometido: cuando toda la carga de la distribución es del
mismo signo. En ese caso el término monopolar no se anula y es el dominante, y medimos el peso de cada término respecto de él.

\begin{teorema}[Estimación del error]\label{teo:cota-error}
Sea $\rho$ una distribución de carga \emph{de un solo signo} (por ejemplo, $\rho\ge0$) con carga neta $Q\neq0$, contenida en la bola de radio $a$ centrada en el origen
($a$ es el \emph{tamaño característico} de la distribución). Entonces, para todo $r>a$ y todo $n\in\N$,
\begin{equation}\label{eq:cota-error}
  \boxed{\;\left\lvert\frac{\phimul{n}(\vc{r})}{\phimul{0}(\vc{r})}\right\rvert\le\Bigl(\frac{a}{r}\Bigr)^{n}\;}
\end{equation}
\end{teorema}
\begin{proof}
Por \cref{eq:desarrollo-multipolar} y \cref{eq:monopolar},
\[
  \frac{\phimul{n}(\vc{r})}{\phimul{0}(\vc{r})}=\frac{\dfrac{1}{4\pi\eo}\dfrac{1}{r^{n+1}}\displaystyle\int_{V'}r'^{\,n}\Pl{n}(\cosp)\,\rho(\rp)\,\dVp}{\dfrac{1}{4\pi\eo}\dfrac{Q}{r}}
  =\frac{1}{Q\,r^{n}}\int_{V'}r'^{\,n}\Pl{n}(\cosp)\,\rho(\rp)\,\dVp .
\]
Como $\rho\ge0$, $\abs{\rho}=\rho$; además $\abs{\Pl{n}(\cosp)}\le1$ (\cref{lem:legendre-cota}) y $r'^{\,n}\le a^{n}$ en el soporte de $\rho$. Por tanto,
\[
  \left\lvert\int_{V'}r'^{\,n}\Pl{n}(\cosp)\,\rho\,\dVp\right\rvert\le\int_{V'}r'^{\,n}\,\abs{\Pl{n}(\cosp)}\,\rho\,\dVp\le a^{n}\int_{V'}\rho\,\dVp=Q\,a^{n},
\]
y dividiendo por $Q\,r^{n}>0$ se obtiene \cref{eq:cota-error}. Si $\rho\le0$ se aplica lo anterior a $-\rho$.
\end{proof}

\begin{corolario}[Error al truncar en el orden $N$]\label{cor:resto}
En las condiciones del \cref{teo:cota-error}, el error relativo del desarrollo truncado (respecto del término dominante) satisface
\begin{equation}\label{eq:resto}
  \frac{\abs{\phi-\phi^{(N)}}}{\abs{\phimul{0}}}\le\sum_{n=N+1}^{\infty}\Bigl(\frac{a}{r}\Bigr)^{n}=\frac{(a/r)^{N+1}}{1-a/r}.
\end{equation}
\end{corolario}
\begin{proof}
$\abs{\phi-\phi^{(N)}}=\bigl\lvert\sum_{n>N}\phimul{n}\bigr\rvert\le\sum_{n>N}\abs{\phimul{n}}\le\abs{\phimul{0}}\sum_{n>N}(a/r)^{n}$, y la serie geométrica de razón $a/r<1$ se suma
directamente.
\end{proof}

\begin{observacion}[Uso práctico: cuántos términos tomar]\label{obs:cuantos}
El criterio del formulario~\cite{FORMULARIO} identifica el error con el peso del término $n$-ésimo, $\phimul{n}/\phimul{0}<(a/r)^{n}$. Si, por ejemplo, $r\sim10a$ y
queremos un error inferior al $1\%$,
\[
  \Bigl(\frac{a}{r}\Bigr)^{n}=\Bigl(\frac{a}{10a}\Bigr)^{n}=10^{-n}\le0{,}01\quad\Longrightarrow\quad n=2,
\]
y habrá que tomar las contribuciones $n=0,1,2$: monopolar, dipolar y cuadrupolar. El \cref{cor:resto} lo confirma con todo rigor: truncando en $N=2$, el error es menor
que $(0{,}1)^{3}/0{,}9\simeq0{,}11\%$, mientras que truncando en $N=1$ solo se garantiza un $1{,}1\%$. En la práctica se usa también la estimación con el tamaño
\emph{medio} de la distribución, $\langle r'\rangle$, en lugar del máximo $a$:
\begin{equation}\label{eq:estimacion-media}
  \frac{\phimul{n}}{\phimul{0}}\lesssim\frac{\langle r'\rangle^{n}}{r^{n}}\,\frac{1}{Q}\int_{V'}\rho\,\dVp=\Bigl(\frac{\langle r'\rangle}{r}\Bigr)^{n},
\end{equation}
que es menos rigurosa pero más ajustada: es la que usaremos en los seminarios.
\end{observacion}

\begin{observacion}[Alcance de la cota]\label{obs:alcance}
\begin{enumerate}[label=(\alph*)]
  \item \emph{La cota es óptima.} Para la carga puntual del \cref{ej:carga-fuera}, sobre el eje ($\theta=0$) se tiene $\phimul{n}/\phimul{0}=(a/r)^{n}\Pl{n}(1)=(a/r)^{n}$: la
        desigualdad \cref{eq:cota-error} se alcanza.
  \item \emph{La hipótesis de un solo signo es esencial.} Si hay cargas de ambos signos, $Q$ puede ser pequeña o nula (un dipolo, un cuadrupolo\dots) y entonces el término
        monopolar deja de ser el dominante: no tiene sentido comparar con él. En ese caso el error debe estimarse respecto del primer término no nulo, como hicimos al comparar el
        dipolo físico con el puntual (\cref{prop:discrepancia}).
\end{enumerate}
\end{observacion}
